\chapter{Introduction}%
\label{chp:introduction}

\section{Motivation}
\label{sec:motivation}

In recent years, the use of \gls{ai} in education, particularly 
in language learning, has been increasingly growing~\cite{d4db0c16}. 
The increase in \gls{ai}-based features in language learning 
platforms such as Duolingo\footnote{ \label{duolingo}
\url{https://blog.duolingo.com/duolingo-max/}}, 
and recently the use of large language models, shows the 
growing recognition of \gls{ai}'s potential to enhance language 
education. These tools offer many advantages, including 
personalized learning experiences for the individual student's 
needs, as well as on-demand access to language support beyond 
the classroom, making them increasingly valuable in modern 
educational environments. Conversational interfaces and 
chatbots, in particular, are effective tools for practising 
a language as they enable learners to engage in interactive 
dialogue and provide instant feedback, simulating the benefits 
of a human tutor, and make learning more accessible for new 
students. The use of these tools in helping with the study 
of widely-spoken modern languages demonstrates how useful 
conversational \gls{ai} in language education can be.

Despite these advances, classical languages such as Latin do 
not benefit from the same amount of digital resources and 
\gls{ai}-supported educational tools that several modern languages 
do, creating a significant gap. Latin learners face several 
unique challenges. Unlike modern languages, Latin is a 
primarily written and academic language, with no living 
community of native speakers. This means that students cannot 
rely on conversational immersion or interaction with fluent 
speakers as language practice. Although there are digital 
resources for Latin, they tend to focus on functionalities 
such as dictionary lookup or morphological parsing. Tools 
such as Whitaker's Words \cite{words}, while valuable for 
their morphological parsing capabilities, offer only lookup 
functionality and do not provide interactive explanations or 
assist learners in a conversational way through natural-language 
questions. Similarly, traditional Latin dictionaries and 
grammar references, although reliable, require learners to 
independently navigate complex grammatical concepts without 
interactive support.

Although Latin has no native speakers today, it is still widely 
studied. It is still relevant in fields such as classical 
studies, history, law, medicine, linguistics, and theology~\cite{clackson2011companion}. 
The existence of a modern student base with limited tools 
introduces the need for a tool that can make Latin more 
accessible to learners, which can be achieved by expanding 
the use of \gls{ai} for this purpose.

Therefore, the motivation of this project is to create a 
conversational \gls{ai} tool that helps with Latin grammar as a 
way to support Latin learners, addressing the gap between 
classical studies and \gls{ai} for language learning.

\section{Problem Definition}%
\label{sec:problem_definition}

Having established the value of \gls{ai} and conversational tools 
in language learning and the gap in such resources for Latin, 
the main problem this project addresses is the absence of a 
conversational \gls{ai} system that can understand and respond to 
Latin grammar questions.

As a result, students studying Latin lack a system that can 
support them when they encounter specific grammatical 
difficulties, and they do not have a conversational tool to 
ask for immediate assistance. Existing computational resources 
for Latin function as dictionaries by providing lookup data, 
but they lack the interactive capability to answer specific 
questions.

Therefore, the specific problem this project seeks to solve 
can be defined as the need for a \gls{ai}-based system capable of 
understanding user intent in natural-language grammar queries. 
The system must accurately identify what grammatical aspect a 
student is asking about from naturally-phrased grammar 
questions. Unlike simple word-based lookup, learners may 
phrase questions in various forms, and the system must be 
able to interpret these and map the questions to the 
grammatical concepts being queried. Additionally, the system 
must also be capable of producing accurate and useful feedback. 
The system must correctly identify and explain the grammatical 
features of words. Furthermore, the system must handle 
ambiguous forms where a single word form may correspond 
to multiple morphological analyses depending on context. 
This information should be delivered conversationally as 
accessible explanations that help learners understand these 
grammatical concepts.

\section{Aims and Objectives}
\label{sec:aims_and_objectives}

The aim of this project is to develop a conversational tool 
for Latin language learners to help them with and provide 
answers to queries related to Latin grammar.

In order to achieve this aim, the following objectives have 
been identified:

% the objectives could be expressed in a measurable way 
% to make it easier to assess the extent to which 
% they are achieved.

\begin{enumerate}
    \item 
\end{enumerate}


%  Construct a dataset of Latin words and their grammatical features by using existing Latin language
% datasets and expanding them.
% • Design and generate a synthetic dataset for intent classification, consisting of grammar-related
% questions covering different grammatical questions.
% This will serve as training data for the intent classification model.
% • Fine-tune a BERT-based model for intent classification to identify user grammatical queries.
% • Develop a Response Generation System that, according to the intent classification, will carry out
% the necessary steps to construct the appropriate
% answer.
% • Implement a conversational interface that accepts
% natural language input in English, so that users
% can interact with and submit grammar questions to
% the model, and receive the output in the form of a
% response.
% • Evaluate the system’s performance through quantitative measures and qualitative user testing with
% Latin learners.


\section{Proposed Solution}
\label{sec:proposed_solution}

\section{Document Structure}
\label{sec:document_structure}

The subsequent chapters of this document are organized as follows:

\begin{itemize}
    \item Chapter~\ref{chp:background}
    \item Chapter~\ref{chp:literature_review} provides a review of the existing literature and related work in the areas relevant to this project.
    \item Chapter~\ref{chp:methodology} details the design and implementation of the system, including the architecture, dataset creation, intent classification model, fine-tuning of Latin BERT for morphological analysis, and response generation.
    \item Chapter~\ref{chp:evaluation} presents and discusses the results obtained from the evaluation and experiments carried out.
    \item Chapter~\ref{chp:conclusion} summarizes the work carried out in this project, discusses its limitations, and suggests areas for future work.
\end{itemize}
